\documentclass[11pt]{article}
\usepackage[margin=1in]{geometry}
\usepackage{amsmath,amssymb,amsthm,hyperref}
\hypersetup{colorlinks=true,urlcolor=blue,linkcolor=blue}
\newtheorem{theorem}{Theorem}
\newtheorem{proposition}{Proposition}
\title{Finite positivity forces binomial rigidity\\\large MathDB 325243}
\author{Principia Math In Progress}
\date{18 September 2026}
\begin{document}
\maketitle

\section*{Task 1: complete classification for every positive integer $s$}
Let $s\ge1$ be an integer, $q\in\mathbb R\setminus\{0\}$, and
$t=\widetilde q>0$. For real $p_0,\ldots,p_s$, define
\begin{equation}\label{eq:g}
 g_k=\sum_{m=0}^{\min(k,s)}
 \binom{s-1+k-m}{s-1}p_mq^m(-t)^{k-m}\qquad(k\ge0).
\end{equation}
L\'oczi and Ketcheson's Conjecture 2.11 in the preprint, numbered
8.2 in the published paper \cite{LK}, asks whether
$g_0,\ldots,g_{s+1}\ge0$ force the binomial coefficients below.
We prove this for every $s$, and give the exact interpretation of
the extra variable $p_{s+1}$ printed in that conjecture.

\begin{theorem}\label{thm:main}
The following conditions are equivalent:
\begin{enumerate}
\item $g_k\ge0$ for $0\le k\le s+1$;
\item for some $c\ge0$,
\[
             p_m=c\binom{s}{m}\left(\frac tq\right)^m
                    \qquad(0\le m\le s).
\]
\end{enumerate}
In either case $c=p_0$, $g_0=c$, and $g_k=0$ for every $k\ge1$.
If $p_{s+1}$ is also supplied as an independent real variable,
it is arbitrary: the complete feasible set is the displayed ray
in $(p_0,\ldots,p_s)$ times $\mathbb R$ in the last coordinate.
\end{theorem}

\begin{proof}
All identities can be read in the formal power series ring
$\mathbb R[[z]]$; no convergence assertion is needed. Set
\[
 P(z)=\sum_{m=0}^s p_mq^m z^m,
 \qquad G(z)=\frac{P(z)}{(1+tz)^s}.
\]
The negative binomial expansion
\[
 (1+tz)^{-s}=\sum_{j\ge0}(-1)^j\binom{s+j-1}{s-1}t^jz^j
\]
shows that $G(z)=\sum_{k\ge0}g_kz^k$, with precisely the sums
in \eqref{eq:g}. Since $\deg P\le s$, the coefficient of $z^{s+1}$
in $(1+tz)^sG(z)=P(z)$ is zero. Thus
\begin{equation}\label{eq:certificate}
      0=\sum_{j=0}^s\binom{s}{j}t^j g_{s+1-j}.
\end{equation}
Every weight $\binom{s}{j}t^j$ is strictly positive.
Under condition (1), every summand is nonnegative, so
\[
                  g_1=g_2=\cdots=g_{s+1}=0.
\]
Also $g_0=p_0\ge0$. Comparing coefficients of degrees $0$ through
$s$ in $(1+tz)^sG=P$ now gives
\[
                 p_mq^m=g_0\binom{s}{m}t^m
                        \qquad(0\le m\le s).
\]
Division by $q^m$ proves (2), including negative $q$.
Conversely, (2) gives $P(z)=c(1+tz)^s$, so $G(z)=c$ and all
the asserted values of $g_k$ follow. Finally, $p_{s+1}$ occurs
in none of the defining sums, so it is unrestricted.
\end{proof}

\paragraph{The printed last index.}
Both source versions also assert the formula for $m=s+1$,
which would force $p_{s+1}=0$. That extra conclusion is false
when $p_{s+1}$ is independent: take $s=q=t=1$ and
$(p_0,p_1,p_2)=(1,1,1)$. Then $(g_0,g_1,g_2)=(1,0,0)$.
Theorem~\ref{thm:main} proves the full intended rigidity assertion
for all the coefficients that enter the hypotheses; it also proves
every claimed vanishing sum and its converse. The complete
all-$s$ classification, rather than this unused-coordinate
counterexample alone, is the mathematical result here.
In the motivating degree-$s$ polynomial setting,
$p_{s+1}=P^{(s+1)}(x)/(s+1)!=0$ automatically.

\begin{proposition}[A general denominator and a sharp cutoff]
Let $D(z)=\sum_{j=0}^s d_jz^j$ with every $d_j>0$, and
$P\in\mathbb R[z]$ with $\deg P\le s$.
If the coefficients $g_0,\ldots,g_{s+1}$ of $P/D$ are nonnegative,
then $P=g_0D$. For each such $D$, testing only degrees
$0,\ldots,s$ does not imply this conclusion.
\end{proposition}
\begin{proof}
The coefficient of $z^{s+1}$ gives
$\sum_{j=0}^s d_jg_{s+1-j}=0$. Positivity again forces
$g_1=\cdots=g_{s+1}=0$, and comparison through degree $s$
gives $P=g_0D$.
For sharpness, fix $\varepsilon>0$ and put
\[
 P(z)=D(z)(1+\varepsilon z)-\varepsilon d_sz^{s+1}.
\]
This has degree at most $s$, and
\[
 \frac{P(z)}{D(z)}=1+\varepsilon z
           -\frac{\varepsilon d_sz^{s+1}}{D(z)}.
\]
Hence $g_0=1$, $g_1=\varepsilon$, and $g_2=\cdots=g_s=0$,
where the last range is empty if $s=1$. But
$g_{s+1}=-\varepsilon d_s/d_0<0$, and $P/D$ is not constant.
\end{proof}

\paragraph{Verification and scope.}
The exact verifier checks the positive linear certificate by
coefficient cancellation, compares direct sums with formal series
division, and checks the feasible family and the sharp-cutoff
construction. The proof for arbitrary $s$ is the coefficient
argument above, not an extrapolation of these computations.
The source's Runge--Kutta optimal-radius questions are separate:
the sign $t>0$ used here must not be replaced by $t=-a<0$
without a new argument.

\section*{Four proposed follow-up tasks}
\begin{description}
\item[Task 2.] For denominators having zero or negative coefficients,
characterize when some finite number of nonnegative coefficients of
$P/D$, with $\deg P\le\deg D$, forces $P/D$ to be constant.
\item[Task 3.] Determine optimal quantitative bounds on the distance
from $P$ to a multiple of $D$ when the tested coefficients are only
bounded below by small errors; study uniformity as the positive
coefficients of $D$ approach zero.
\item[Task 4.] For $\deg P>\deg D$ and a specified finite testing
range, classify the extreme rays of the cone defined by
coefficient nonnegativity of $P/D$.
\item[Task 5.] Find a replacement for the positive certificate
\eqref{eq:certificate} compatible with the denominator $(1-az)^s$
and the order conditions in the source's general rational-function
optimal-radius conjectures.
\end{description}
These are proposed research directions and receive no additional
solution credit.

\begin{thebibliography}{9}
\bibitem{LK} L. L\'oczi and D. I. Ketcheson,
\emph{Rational functions with maximal radius of absolute monotonicity},
LMS Journal of Computation and Mathematics \textbf{17}(1) (2014),
159--205. Conjecture 8.2, p.~197.
\href{https://doi.org/10.1112/S1461157013000326}{doi:10.1112/S1461157013000326}.
Preprint: \href{https://arxiv.org/abs/1303.6651}{arXiv:1303.6651v1},
Conjecture 2.11, p.~10.
\end{thebibliography}
\end{document}
